

\section{\ours: \textbf{Vi}suo-\textbf{Ta}ctile \textbf{L}atent Steering}
\label{sec:method}
\ours addresses the challenges in Sec.~\ref{sec:open_challenges} by decomposing multimodal steering into long-horizon visual mode selection and short-horizon tactile refinement. The visual stage selects an action mode based on predicted semantic progress, while the tactile stage refines it using predicted contact quality, matching the complementary timescales of vision and touch. We formulate this bi-level objective in Sec.~\ref{subsec:method_bilevel}, introduce a visuo-tactile latent world model for outcome prediction in Sec.~\ref{subsec:method_world_model}, define visual and tactile verifiers in Sec.~\ref{subsec:method_reward_modeling}, and describe the overall steering approach in Sec.~\ref{subsec:method_bilevel_execution}.















\subsection{Bi-level Visuo-Tactile Steering Formulation}
\label{subsec:method_bilevel}
We formulate the multimodal steering problem as a bi-level optimization problem.
For both levels of the optimization, we use a latent visuo-tactile world model to predict the outcomes of any action sequence, $\acttraj$, across both modalities:
$\latenttrajvpred, \latenttrajtpred \sim  p_\phi\!\left(
\cdot,\cdot
\mid
\latentv, \latentt, \acttraj
\right)$ where $\latent \in \latentSpace$ are the latent states obtained from the respective modality's encoder $\enc^m(\obs), m \in \{v, \tau\}$. 
The \textit{inner} optimization selects a globally promising action mode $\bar{\acttraj}_{t:t+H}$ with visual task reward $R^v(\latenttrajvpred; \task)$. 
The \textit{outer} optimization stays close to, but refines, the first $h$ steps of the visual plan $\bar{\acttraj}_{t:t+h}$ using the tactile latents and reward $R^\tau(\latenttrajtpred; \task)$.
The final action chunk informed by both vision and touch is denoted by $\acttraj^\star_{t:t+h}$:
\begin{equation}
\begin{aligned}
\mathbf{a}^*_{t:t+h}
&=
\arg\max_{\acttraj_{t:t+h}}
\overbrace{
\log p_\theta\!\left(
\acttraj_{t:t+h}
\mid
\textcolor{ocean}{\bar{\mathbf{a}}_{t:t+h}}, o_t
\right)
}^{\text{visual prior around selected mode}}
+
\beta \cdot
\overbrace{
R^\tau\!\left(
\textcolor{solid_green}{\latenttrajtpred_{t:t+h}}; \mathcal{L}
\right)
}^{\text{\textcolor{solid_green}{\textbf{tactile}}: local contact reward}}
\\
&\text{s.t.}\quad
\textcolor{ocean}{\bar{\mathbf{a}}_{t:t+H}}
=
\underbrace{
\arg\max_{\mathbf{a}_{t:t+H}\sim \pi_\theta(\cdot\mid o_t)}
R^v\!\left(
\textcolor{ocean}{\latenttrajvpred_{t:t+H}}; \mathcal{L}
\right)
}_{\text{\textcolor{ocean}{\textbf{vision}:} select action mode via visual latents \& reward}},
\\
& \qquad \textcolor{ocean}{\latenttrajvpred_{t:t+h'}}, \textcolor{solid_green}{\latenttrajtpred_{t:t+h'}}
\sim
\underbrace{p_\phi\!\left(
\cdot,\cdot
\mid
\latentv_t, \latentt_t, \acttraj_{t:t+h'}
\right)}_{\text{latent visuo-tactile world model}},
\quad h' \in \{h, H\}, \\ 
&\qquad \latentv_t =\encv(\obsv_t), \quad \latentt_t =\enct(\obst_t).
\end{aligned}
\label{eq:bilevel_joint_annotated}
\end{equation}







In the \textit{inner} optimization, actions are sampled from the base policy, $\acttraj_{t:t+H} \sim \pi_\theta(\cdot \mid \obs_t)$, rolled out with the latent visuo-tactile dynamics model $p_\phi$, and the actions with the highest long-horizon visual reward $R^v(\latenttrajvpred_{t:t+H}; \task)$ are chosen to encourage global task progress. 
The \textit{outer} optimization uses the visually-selected action mode, $\bar{\acttraj}$, as an initial behavior proposal. 
The first term in the optimization, $p_\theta$, is a visual action prior (described in detail in Sec.~\ref{subsec:method_bilevel_execution}) centered on this action mode; this encourages the optimization to remain close to the globally desired behavior. 
The second term locally refines this visually selected mode by reweighting shorter $h$-step action chunks according to the tactile reward $R^\tau(\latenttrajtpred_{t:t+h}; \task)$. 
The hyperparameter $\beta$ controls the strength of tactile refinement relative to the visual prior. 
Intuitively, our formulation uses vision for global mode selection and then touch for local contact refinement. 






    
\subsection{Visuo-Tactile Latent World Model for Multi-Modal Latent Outcome Prediction}
\label{subsec:method_world_model}

Eq.~\ref{eq:bilevel_joint_annotated} uses a latent world model $p_\phi$ to predict the visual and tactile futures induced by candidate action sequences. 
We encode the current visual and tactile observations separately as
$
    \latentv_t = \encv(\obsv_t), \;
    \latentt_t = \enct(\obst_t), \;
$
with joint vision-tactile state $\latent_t = (\latentv_t, \latentt_t)$. In our implementation, $\encv$ is the pretrained DINOv3 encoder~\cite{simeoni2025dinov3} and $\enct$ is the pretrained AnyTouch2 encoder~\cite{fenganytouch2} with structured latent spaces for scene understanding and tactile reasoning.
On top of the frozen encoders, we train an action-conditioned latent dynamics model inspired by~\cite{zhoudino}. Starting from $\hat{\latent}_t=\latent_t$, the model recursively predicts future joint visuo-tactile latents over action chunks:
$    \hat{\latent}_{t+h}
    =
    p_\phi
    \left(
        \hat{\latent}_{t},
        \action_{t:t+h}
    \right).$
We train the model 
with a multi-step latent prediction objective following~\cite{higuera2026visuo} with details in App~\ref{app:implement_visuo_tactile_world_model}.
We additionally train decoders
$
    \hat{\obs}^v = \decv(\latentv), \;
    \hat{\obs}^\tau = \dect(\latentt)
$
for visualization.
Together, this multimodal model enables long-horizon visual lookahead as well as short-horizon tactile prediction.


\subsection{Multi-modal Reward Modeling for Verification}
\label{subsec:method_reward_modeling}
After the latent world model predicts multimodal (latent) outcomes, we need to verify how well these predictions align with global task progress and local contact quality, which may vary across task phases. 
We use language-conditioned visual and tactile verifiers, with phase-dependent textual objectives specifying the relevant criteria for each modality.


\para{Phase-Dependent Verification Objectives}
A single task instruction is often too coarse for long-horizon manipulation: e.g., $\task$=``transfer liquid to the blue cup and return'' does not explicitly specify intermediate visual objectives (such as moving back to the red cup) or contact requirements (such as squeezing only during dispensing). We therefore use a VLM offline to decompose the task into phase-level textual objectives,
    $\{\lang_p\}_{p=1}^{P} = g^{\mathrm{VLM}}(\task)$, 
where each phase provides a visual objective $\lang_p^v$ and a tactile objective $\lang_p^\tau$. The visual verifier determines when to switch to the next phase.

\para{Visual Verifier}
Visual success depends on if the imagined outcomes appear to be making progress toward the visual objective described in $\lang_p^v$. 
Recent work has developed strong, off-the-shelf, general-purpose reward models that take as input RGB images and textual task descriptions \cite{tan2025robo, liang2026robometer, ma2025vision}.
We leverage these models zero-shot, and we decode our predicted visual latents into RGB image observations $\hat{\obs}^v_{t+H} = \decv(\hat{\latent}^v_{t+H})$ and concatenate that with the previous history of observations we have obtained $\obs_{\leq t}$ as the input of ROBOMETER~\cite{liang2026robometer} reward model for online visual evaluation. 
This reward is used for long-horizon visual mode selection, where the goal is to distinguish globally meaningful task progress, like moving to the correct cup, mark, or hole.

\para{Tactile Verifier}
Tactile success depends on fine-grained, often visually ambiguous interactions, such as stable grasping versus incipient slip. 
We introduce a tactile verifier that is \textit{discriminative} over subtle contact outcomes, \textit{efficient} for inference-time  evaluation, and \textit{flexible} to adapt to phase-specific textual contact goals.
Specifically, we use a pretrained tactile encoder~\cite{fenganytouch2} with a semantically-aligned latent space and define the tactile reward as the cosine similarity between the phase-level contact description and the tactile embedding: 
$$R^\tau(\hat{\latent}^\tau_{t+h}, \lang_p^\tau)
=
\cos\left(
    \hat{\latent}^\tau_{t+h},
    \enc^\lang(\lang_p^\tau)
\right),$$
where $\enc^\lang(\cdot)$ is the corresponding text encoder. Computing rewards directly in this latent space makes the verifier discriminative over subtle contact states, efficient for steering without tactile image decoding, and flexible to language-specified contact goals such as ``grasp lightly'' or ``grasp heavily.'' To the best of our knowledge, this is the first general-purpose language-conditioned tactile reward. 






\subsection{Inference-time Optimization for Bi-level Visuo-Tactile Steering Objective }
\label{subsec:method_bilevel_execution}
At inference time, we solve the bi-level optimization problem from Sec.~\ref{subsec:method_bilevel} by combining sampling-based visual mode selection with tactile-guided diffusion editing. We describe the implementation details behind this optimization in this section, with further details in Algo.~\ref{alg:vital} in App.~\ref{app:implement_visual_steering}.

\para{Sampling \& Verification as Visual Steering}
As shown in the inner optimization in Eq.~\ref{eq:bilevel_joint_annotated}, the visual stage selects a global behavior mode by evaluating long-horizon candidate action sequences. We obtain a set of long-horizon action sequences
$
\mathbb{A}_N=\{\acttraj_{t:t+H}^{(i)} \mid 
\acttraj_{t:t+H}^{(i)} \sim \pi_\theta(\cdot \mid \obs_t)\}_{i=1}^{N}
$
by rolling-out the short action samples obtained from the policy within the world model, decoding the visual observations, passing these into the policy, and repeating. We choose the sequence with the highest visual reward:
$
\bar{\acttraj}_{t:t+H}
=
\arg\max_{\acttraj_{t:t+H} \in \mathbb{A}_N}
R^v
\left(
    \hat{\latenttraj}^v_{t:t+H}; \lang_p^v
\right).$
We use sampling-based optimization here because of its simplicity and stability; alternatives like classifier-based guidance would  
require differentiating through  through the VLM, latent world model, and recursive policy generation, resulting in potentially unstable gradients.  
Thus, we use the visual verifier to rank candidate futures rather than guide denoising directly. 
The selected sequence $\bar{\action}_{t:t+H}$ serves as the visual anchor, and its first $h$ steps are passed to the outer optimization for tactile refinement.





\para{Diffusion Editing as Tactile Steering}
Conditioned on the visual anchor $\bar{\acttraj}_t :=\bar{\acttraj}_{t:t+h}$, the goal of tactile steering is to locally refine the action chunk to satisfy contact requirements while preserving the action mode selected by the visual steering. We evaluate tactile refinement using the short-horizon tactile reward in Sec.~\ref{subsec:method_reward_modeling}. The short horizon $h$ is important because tactile signals are most informative for immediate contact correction, such as adjusting force or stabilizing a grasp.


Following Eq.~\ref{eq:bilevel_joint_annotated} and prior diffusion editing methods~\cite{meng2021sdedit}, we define a reference-conditioned prior around visual anchor $\bar{\acttraj}_{t}$:
$
p_\theta(\acttraj_{t} \mid \bar{\acttraj}_{t}, \obs_t)
:=
\int
    \rho_\theta(\acttraj_{t} \mid \mathbf{x}_K, \obs_t)
\nu_K(\mathbf{x}_K \mid \bar{\acttraj}_t)
d\mathbf{x}_K,
$
where $\nu_K$ partially noises $\bar{\acttraj}_t$ to diffusion level $K$ and $\rho_\theta$ denoises it under the base policy. This prior captures plausible local edits of the visual anchor, with smaller $K$ enforcing closer refinements and larger $K$ allowing larger deviation. We then define the tactile-refined policy $p^\textrm{refine}$ by reweighting the policy distribution with predicted tactile reward,
$
\log p^{\mathrm{refine}}
=
\log p_\theta(\acttraj_t \mid \bar{\acttraj}_t, \obs_t)
+
\beta R^\tau(\hat{z}^\tau_{t+h}, \lang^\tau_p),
$
and approximate sampling from $p^\textrm{refine}$ via tactile-guided reverse diffusion: we partially noise $\bar{\acttraj}_t$ and bias denoising toward actions with higher tactile reward. Thus, the visual stage selects \emph{what} global behavior to execute, while the tactile stage refines \emph{how} to execute it with touch. Details are in App.~\ref{app:implement_visual_steering}-~\ref{app:implement_tactile_steering}.



